\subsection*{4.2 Singular and Nonsingular variables}
Let us fix the finite field $\mathbb F_p$, $p>2$ for the rest of the article. We also fix a set of unsatisfiable polynomials $F$ over Boolean variables $\mathbf x\cup\mathbf y$, and a set of extension axioms $\mathsf{Ext}$ of the form $z-Q$ over variables $\mathbf z$. Whenever we refer to a refutation $\Pi$, we assume that it is a PC+Ext refutation of $F\cup\mathsf{Ext}$.
\paragraph{Definition10 (Support of a variable).}
Let $z-Q(w_{i_1},\ldots,w_{i_{\kappa}})=0$ be a $\kappa$-local extension axiom associated with $z$. We define the set $\vars(Q)=\{w_{i_1},\ldots,w_{i_{\kappa}}\}$ and sometimes write $\vars(z)$ to denote $\vars(Q)$, the set of variables that $z$ depends on. The support of $z$, $\supp(z)\subseteq[0,p-1]$, is equal to the set of all values $a\in[0,p-1]$ such that there exists a Boolean assignment $\alpha$ to the variables of $Q$ such that $Q(\alpha)=a$. Sometimes we also indicate this by $\supp(Q)$.
We extend the definition of support also to Boolean variables. For a Boolean variable $w$, $\supp(w)=\{0,1\}$ as enforced by the Boolean axiom $w^2=w$.
\paragraph{Definition11 (Singular and Nonsingular variables w.r.t. Ext).}
Let $\mathsf{Ext}$ be a set of extension axioms and let $z$ be an extension variable with an axiom in $\mathsf{Ext}$. We say that $z$ is Singular w.r.t. $\mathsf{Ext}$ iff $0\in\supp(z)$; otherwise we say that $z$ is Nonsingular w.r.t. $\mathsf{Ext}$. Any Boolean variable is considered Singular by default, independent of the set $\mathsf{Ext}$, since zero always belongs to its support. For a term $t$, let $\sing(t)$ be the subterm of $t$ containing the Singular variables in $t$, and let $\nsing(t)$ be the subterm of $t$ containing the Nonsingular variables.
Note that for a Singular extension variable $z$, it is possible to set $z$ to zero, However, we note that this may falsify other polynomial equations in $F$. For example, if $xy=0$ is a polynomial equation in $F$, and the extension axiom for $z$ is $z-1+xy=0$, then setting $x=y=1$ forces $z=0$, but this falsifies $xy=0$.
\paragraph{Definition12.} Let $A\subseteq[1,\ldots,p-1]$, $A\neq\emptyset$. Define $\ell(A)$ to be the least $\ell\in[1,p-1]$ such that the set $\{a^{\ell}\mid a\in A\}$ is singleton. For a Nonsingular $z$, define $\ell(z)=\ell(\supp(z))$.
\paragraph{Lemma13.} Let $z$ be a Nonsingular extension variable with extension axiom $z-Q=0$. Then the following polynomial equations are implied by (and therefore derivable from) the extension axiom for $z$ plus the Boolean axioms for all variables in $\vars(Q)$, in degree at most $|\vars(Q)|$.
\begin{enumerate}
\item $z-Q'=0$, where $Q'$ is the multilinear version of $Q$;
\item For any $A'\subseteq[0,p-1]$ such that $\supp(z)\subseteq A'$, $\prod_{a\in A'}(z-a)=0$;
\item $z^{\ell(z)}-c=0$ for some $c\in\mathbb F_p^*$.
\end{enumerate}
In particular, if $z$ is Nonsingular, then the polynomial equation $z^{p-1}-1=0$ is implied by $z-Q=0$ together with the Boolean axioms for $\vars(Q)$.
\begin{proof}
Let $z-Q(w_{i_1},\ldots,w_{i_{\kappa}})=0$ be the extension axiom for $z$, and let $\supp(z)=A\subseteq A'\subseteq[1,p-1]$. First, we can derive the multilinear version of $Q$, $Q'$, from $Q$ together with the Boolean axoms $w^2-w=0$ for all $w\in\vars(Q)$. Secondly, by definition, $\supp(z)=A$ means that the allowable values for $z$ over Boolean assignments to $\vars(Q)$ are the values in $A$. Therefore, $z-Q=0$ together with the Boolean axioms $w^2-w=0$ for all $w\in\vars(Q)$ implies $\prod_{a\in A}(z-a)=0$. Furthermore, this polynomial has a PC derivation, by the derivational completeness of PC. Since $A\subseteq A'$, $\prod_{a\in A'}(z-a)=0$ is a weakening of $\prod_{a\in A}(z-a)=0$ and is therefore derivable from $\prod_{a\in A}(z=a)=0$. Lastly, we will argue that there exists some constant $c\in\mathbb F_p^*$ such that $z^{\ell(A)}-c=0$ is semantically implied by $z-Q=0$ plus the Boolean axioms for $\vars(z)$ and therefore is derivable from these axioms. Since the only allowable values for $z$ under the Boolean axioms are the values in $A$, and since by definition of $\ell(A)$, for every $a\in A$, $a^{\ell(A)}=c$ for some $c\in\mathbb F_p^*$, it follows that $z^{\ell(A)}-c=0$.
\end{proof}
\paragraph{Definition14.} For a term $t$ and a variable $w$, $\deg(t,w)$ is equal to the degree of $w$ in $t$. If $w$ is Nonsingular, then $w^{p-1}=1\bmod p$, so $\deg(t,w)<p-1$. On the other hand if $w$ is Singular then we have $w^p=w\bmod p$ and therefore $\deg(t,w)<p$. For a term $t$ the degree of $t$, $\deg(t)$, equals $\sum_{w\in\vars(t)}\deg(t,w)$.
\subsection*{4.3 Quadratic degree}
The next definition is a generalization/modification of Sokolov's definition of Quadratic degree for the more general scenario where the proof contains extension variables that are Singular as well as ones that are Nonsingular.
\paragraph{Definition15 (Quadratic degree).}
Let $V$ be a set of variables and let $S$ be a subset of $V$. For a pair of terms $t_1,t_2$ over $V$, and a variable $w\in V$, we define $\Qdeg_S(t_1,t_2,w)$ as follows. If $w\in S$, then $\Qdeg_S(t_1,t_2,w)=1$ if $w$ occurs in at least one of $t_1$ or $t_2$; if $w\notin S$, then $\Qdeg_S(t_1,t_2,w)=1$ if and only if $\deg(t_1,w)\neq\deg(t_2,w)$. The overall quadratic degree of the pair $t_1,t_2$, $\Qdeg_S(t_1,t_2)$, is equal to $\sum_{w\in V}\Qdeg_S(t_1,t_2,w)$. The quadratic degree of a polynomial $P$ is equal to the maximum quadratic degree over all pairs $(t_1,t_2)$ such that $t_1,t_2\in P$. For a proof $\Pi$, the quadratic degree of $\Pi$ is the maximum quadratic degree over all polynomials $P\in\Pi$.
We usually instantiate the above definition with $V=\mathbf x\cup\mathbf y\cup\mathbf z$ and with $S$ being the set of Singular variables as defined by the extension axioms corresponding to $\mathbf z$. However since $\Qdeg_S$ is a different measure for every $S$, and our set of Singular variables can change under the application of a restriction $\rho$ to the variables in $V$, we must make sure that our measure of Quadratic degree does not change significantly under a restriction.\footnote{If the set $S$ does not change under a restriction, $\Qdeg_S$ can still change under the restriction as terms can shrink or disappear when variables are set by the restriction. However, this is no different from how the usual notion of degree changes under a restriction, and it is trivial to show that $\Qdeg_S$ always decreases. Therefore we ignore this for the sake of simplicity.} Fortunately, we can show that for any two sets $S$ and $T$ such that $T\subseteq S$, $\Qdeg_T\leq\Qdeg_S$. Along with the simple observation that the set of Singular extension variables can only shrink under a restriction, this implies that our measure of Quadratic degree can only decrease under a restriction. We make this formal below.
\paragraph{Lemma16.} Let $V$ be a set of variables and let $S$ and $T$ be subsets of $V$ such that $T\subseteq S$. Then for any two terms $t_1,t_2$ over $V$, $\Qdeg_T(t_1,t_2)\leq\Qdeg_S(t_1,t_2)$.
\begin{proof}
Note that for a variable $w\in S-T$, $\Qdeg_S(t_1,t_2,w)=1$ when $w$ has a nonzero exponent in one of $t_1$ or $t_2$, otherwise zero. However, $\Qdeg_T(t_1,t_2,w)=1$ if and only if the previous condition is satisfied and the exponents of $w$ in $t_1$ and $t_2$ are not equal. Thus the claim follows.
\end{proof}
Henceforth, when we refer to Quadratic degree, we always fix the set $S$ to be the set of Singular variables w.r.t. the underlying extension axioms. We have the following important corollary that this measure always decreases under a restriction to the underlying variables.
\paragraph{Corollary17.} Let $F$ be a set of unsatisfiable polynomials over variables $\mathbf x\cup\mathbf y$ and let $\mathsf{Ext}$ be a set of extension axioms of the form $z-Q(w_{i_1},\ldots,w_{i_{\kappa}})$ for variables $z\in\mathbf z$ and $w_{i_1},\ldots,w_{i_{\kappa}}\in\mathbf x\cup\mathbf y$. Let $\rho$ be a restriction to $\mathbf x\cup\mathbf y$ and let $\mathsf{Ext}|_{\rho}$ be the axioms given by $z-Q|_{\rho}$ for each axiom $z-Q\in\mathsf{Ext}$. The Quadratic degree w.r.t. $\mathsf{Ext}|_{\rho}$ is at most the Quadratic degree w.r.t. $\mathsf{Ext}$.
\begin{proof}
Since $\supp(Q|_{\rho})\subseteq\supp(Q)$ for any polynomial $Q$, we have that the set of Singular variables under $\mathsf{Ext}|_{\rho}$ is a subset of those under $\mathsf{Ext}$. Therefore our claim follows from the previous lemma.
\end{proof}
\paragraph{Lemma18 (Quadratic degree upper bounds degree of Singular variables).} For any term $t$, $\deg(\sing(t))\leq p\Qdeg(t,t)$.
\begin{proof}
For any Singular variable $w$, $\Qdeg(t,t,w)=1$ if and only if $w$ occurs in $t$. Since $w$ can occur in $t$ with degree at most $p-1$, the claim follows.
\end{proof}
\paragraph{Definition19 (High quadratic degree terms).} For a proof $\Pi$ and $d\geq0$, let $H_d(\Pi)$ denote the set of unordered pairs $(t_1,t_2)$ of quadratic degree at least $d$. That is, $H_d(\Pi)$ is the set of unordered pairs of terms $(t_1,t_2)$ such that $t_1,t_2$ both occur in $P$ for some polynomial $P\in\Pi$, and $\Qdeg(t_1,t_2)\geq d$.
\paragraph{Lemma20.} Let $\Pi$ be a PC+Ext refutation of $F$ and let $z$ be a Nonsingular variable. Let $\Pi'$ be the proof obtained from $\Pi$ by reducing each line of $\Pi$ by $z^{\ell(z)}-c=0$ for some $c\in\mathbb F_p^*$. Then $|H_d(\Pi')|\leq|H_d(\Pi)|$ for any $d\geq0$.
\begin{proof}
Consider a polynomial $P\in\Pi$ and a pair of terms $(t_1,t_2)$ that occur in $P$. For any variable $w$ distinct from $z$, $\Qdeg(t_1,t_2,w)$ is unaltered when $P$ is reduced by $z^{\ell(z)}=c$. On the other hand, if $z$ does not contribute to the Quadratic degree of $(t_1,t_2)$ i.e. $\Qdeg(t_1,t_2,z)=0$, then it will still be $0$ after reducing by $z^{\ell(z)}=c$. Therefore $\Qdeg(t_1,t_2)$ never increases for any pair $(t_1,t_2)$ and thus $|H_d(\Pi')|\leq|H_d(\Pi)|$.
\end{proof}
The following is a generalized version of the argument from [20] that shows how to convert a proof with low Quadratic degree to one with low degree.
\paragraph{Lemma21.} Let $F$ be a set of unsatisfiable polynomials of degree $d_0$ with a PC refutation of Quadratic degree at most $d\geq d_0$ over $\mathbb F_p$. Then $F$ has a PC refutation of degree at most $3pd$.
\begin{proof}
The proof of this lemma is largely based on (a slightly cleaner version of) Sokolov's argument ([20], Lemma3.6) that low Quadratic degree over $\{\pm1\}$ variables implies low degree. Our first observation is that Sokolov's argument can be applied to any refutation of low Quadratic degree over $\mathbb F_p$ such that every term contains only Nonsingular variables. In particular if $\{P_j\}$ is a refutation that only contains Nonsingular terms, then we can use his argument to show that $\{t_j^{p-2}P_j\}$ is also a valid refutation for some carefully chosen term $t_j\in P_j$. Moreover, the degree of the latter refutation is bounded by a constant times the Quadratic degree of the former one. To see this, first note that for two Nonsingular terms $t_1$ and $t_2$, we have that $\deg(t_1t_2^{p-2})\leq(p-1)\cdot\Qdeg(t_1,t_2)$, because of the following. For a variable $z$ that is Nonsingular such that $z$ occurs in $t_1$ and $t_2$ with $\deg(t_1,z)=\deg(t_2,z)$, we have $\deg(t_1t_2^{p-2},z)=\Qdeg(t_1,t_2,z)=0$ since it would appear in $t_1t_2^{p-2}$ with an exponent that is a multiple of $p-1$, and $z^{p-1}=1$ holds for Nonsingular variables. Any other Nonsingular $z$ that occurs in at least one of $t_1$ and $t_2$ has $\deg(t_1t_2^{p-2},z)<p-1$ and $\Qdeg(t_1,t_2,z)=1$. Therefore the degree of $t_1t_2^{p-2}$ is at most $p\cdot\Qdeg(t_1,t_2)$ when $t_1$ and $t_2$ contain only Nonsingular variables. This implies that the lines in the new refutation $\{t_j^{p-2}P_j\}$ have degree at most $p$ times the Quadratic degree of the original refutation $\{P_j\}$. Sokolov additionally showed that each line in the new refutation can be derived from previous lines without exceeding degree equal to $2p$ times the Quadratic degree of the original refutation, completing the argument.
In our case we deal with terms containing both Singular and Nonsingular variables. The above argument cannot be applied directly to our case, since it crucially depends on the fact that Nonsingular variables can be raised to the power $p-1$ to make them vanish. Fortunately by Lemma18, the degree of Singular variables in any term is at most $p$ times the Quadratic degree with itself. Given this bound, we can ignore for each term the part that contains Singular variables, and apply the above argument only with respect to the Nonsingular part of each term, to reduce the degree of Nonsingular variables in each term of the refutation. Since we now have a bound on the degree of both Singular and Nonsingular variables in each term, we have bounded its degree. We describe this in full technical detail below.
Let $\{P_j\}$ be a refutation of $F$ with Quadratic degree bounded by $d$. For any term $t$ recall that $\nsing(t)$ denotes the subterm of $t$ containing only Nonsingular variables. Note that $\nsing(t)^{p-1}=1$ for any $t$. For every line $P_j$ in the refutation, we pick a term $t_j\in P_j$ and define $P'_j=\nsing(t_j)^{p-2}P_j$. Note that by the arguments outlined above, for any two terms $t_1$ and $t_2$ in $P_j$, we have $\deg(\nsing(t_1)^{p-2}\nsing(t_2))\leq pd$ and thus the degree of Nonsingular variables in any term of $P'_j$ is bounded by $pd$. Since the Singular variables in any term remain unchanged under multiplication by $\nsing(t_j)^{p-2}$, the Singular degree of $P'_j$ the same as that of $P_j$ and is bounded by $pd$ (Lemma18) and therefore $\deg(P'_j)\leq pd+pd=2pd$. We now show that the set $\{P'_j\}$ forms a valid refutation of $F$ and each $P'_j$ can be derived from previous lines in degree $3pd$. If $P_j$ is one of the axioms, we multiply by $\nsing(t_j)^{p-2}$ to get $P'_j$ for an arbitrary $t_j\in P_j$, and this takes degree $pd_0\leq pd$.
If $P_j=wP_{j_1}$ for $j_1<j$ and some variable $w$, we choose $t_j\in P_j$ such that $t_j=wt_{j_1}$ where $t_{j_1}\in P_{j_1}$ was chosen earlier. If $w$ is Singular, we have $\nsing(t_j)=\nsing(t_{j_1})$ and therefore $P'_j=\nsing(t_j)^{p-2}P_j=w\cdot\nsing(t_{j_1})^{p-2}P_{j_1}=wP'_{j_1}$. On the other hand, if $w$ is Nonsingular, we have $\nsing(t_j)=w\cdot\nsing(t_{j_1})$ and therefore $P'_j=\nsing(t_j)^{p-2}P_j=w^{p-1}\cdot\nsing(t_{j_1})^{p-2}P_{j_1}=P'_{j_1}$. Finally, let $P_j=P_{j_1}+P_{j_2}$ for $j_1,j_2<j$. We pick an arbitrary term $t_j\in P_j$. Note that since $\nsing(t)^{p-1}=1$ for any term $t$, $P_{j_1}=\nsing(t_{j_1})P'_{j_1}$ and $P_{j_2}=\nsing(t_{j_2})P'_{j_2}$ and thus we have
\begin{multline*}
P'_j=\nsing(t_j)^{p-2}\nsing(t_{j_1})P'_{j_1}\\
+\nsing(t_j)^{p-2}\nsing(t_{j_2})P'_{j_2}
\end{multline*}
for $t_{j_1}\in P_{j_1}$ and $t_{j_2}\in P_{j_2}$ chosen earlier. We now show that $\deg(\nsing(t_j)^{p-2}\nsing(t_{j_1}))\leq pd$ and $\deg(\nsing(t_j)^{p-2}\nsing(t_{j_2}))\leq pd$ to conclude the proof. Since every term in $P_j$ appears in one of $P_{j_1},P_{j_2}$, let $t_j\in P_{j_1}$ without loss of generality. Then we have that $t_j,t_{j_1}$ both appear in $P_{j_1}$ and thus $\deg(\nsing(t_j)^{p-2}\nsing(t_{j_1}))\leq pd$. If $t_{j_2}\in P_j$ i.e. it is not cancelled in the sum $P_{j_1}+P_{j_2}$, then we have $t_j,t_{j_2}$ both appear in $P_j$ and hence $\deg(\nsing(t_j)^{p-2}\nsing(t_{j_2}))\leq pd$. If $t_{j_2}\notin P_j$, this implies that it was cancelled in the sum $P_{j_1}+P_{j_2}$ and therefore $t_{j_2}\in P_{j_1}$ and $\deg(\nsing(t_j)^{p-2}\nsing(t_{j_2}))\leq pd$.
\end{proof}
